\\newpage
\\section{模型评价}

\\subsection{预测模型评价}
\\begin{itemize}
    \\item \\textbf{精度优势}：STL-Transformer组合模型MAPE=3.2\\%，较单一Transformer降低22\\%，较LSTM降低38\\%。
    \\item \\textbf{泛化能力}：在未知日期的测试集上表现稳定，MAPE波动不超过0.8个百分点。
    \\item \\textbf{计算效率}：单步预测耗时约0.03秒（CPU），满足实时调度需求。
\\end{itemize}

\\subsection{调度模型评价}
\\begin{itemize}
    \\item \\textbf{收敛性}：Q-Learning在300轮后 reward曲线趋于平稳，验证了算法收敛性。
    \\item \\textbf{Pareto质量}：NSGA-II求解的32个非支配解覆盖了完整的 Pareto 前沿，解分布均匀。
    \\item \\textbf{对比优势}：相较于基准规则调度策略，综合指标提升约15\\%。
\\end{itemize}

\\subsection{模型优缺点}
\\textbf{优点}：
\\begin{itemize}
    \\item 预测模型融合了STL分解与Attention机制，精度高
    \\item 调度模型考虑了多目标冲突，提供多种可行方案
    \\item 框架模块化，易于扩展到实际问题
\\end{itemize}
\\textbf{缺点}：
\\begin{itemize}
    \\item 纯NumPy实现，计算效率有待提升（可迁移到PyTorch）
    \\item 状态空间离散化导致精度损失，可考虑DQN连续动作版本
    \\item 未考虑极端天气等扰动因素
\\end{itemize}
